\documentclass[12pt,a4paper]{article}
\usepackage{fontawesome5}
\usepackage{xcolor}

\usepackage[T1]{fontenc}
\usepackage[utf8]{inputenc}
\usepackage{dsfont} 
\usepackage[polish]{babel}
\usepackage{amsmath}
\usepackage{graphicx}
\usepackage[top=1in, bottom=1.5in, left=1.25in, right=1.25in]{geometry}

\usepackage{subfig}
\usepackage{multirow}
\usepackage{multicol}
\graphicspath{{Images/}}
\usepackage{xcolor,colortbl}
\usepackage{float}
\usepackage{hyperref}
\usepackage{listings}

\newcommand \comment[1]{\textbf{\textcolor{red}{#1}}}

%\usepackage{float}
\usepackage{fancyhdr} % Required for custom headers
\usepackage{lastpage} % Required to determine the last page for the footer
\usepackage{extramarks} % Required for headers and footers
\usepackage{indentfirst}
\usepackage{placeins}
\usepackage{scalefnt}
\usepackage{xcolor,listings}
\usepackage{textcomp}
\usepackage{color}
\usepackage{verbatim}
\usepackage{framed}


\definecolor{codegreen}{rgb}{0,0.6,0}
\definecolor{codegray}{rgb}{0.5,0.5,0.5}
\definecolor{codepurple}{HTML}{C42043}
\definecolor{backcolour}{HTML}{F2F2F2}
\definecolor{bookColor}{cmyk}{0,0,0,0.90}  
\color{bookColor}

\lstset{upquote=true}

\lstdefinestyle{mystyle}{
	backgroundcolor=\color{backcolour},   
	commentstyle=\color{codegreen},
	keywordstyle=\color{codepurple},
	numberstyle=\numberstyle,
	stringstyle=\color{codepurple},
	basicstyle=\footnotesize\ttfamily,
	breakatwhitespace=false,
	breaklines=true,
	captionpos=b,
	keepspaces=true,
	numbers=left,
	numbersep=10pt,
	showspaces=false,
	showstringspaces=false,
	showtabs=false,
}
\lstset{style=mystyle}

\newcommand\numberstyle[1]{%
	\footnotesize
	\color{codegray}%
	\ttfamily
	\ifnum#1<10 0\fi#1 |%
}

\definecolor{shadecolor}{HTML}{F2F2F2}

\newenvironment{sqltable}%
{\snugshade\verbatim}%
{\endverbatim\endsnugshade}

% Margins
\addtolength{\footskip}{0cm}
\addtolength{\textwidth}{1.4cm}
\addtolength{\oddsidemargin}{-.7cm}

\addtolength{\textheight}{1.6cm}
%\addtolength{\topmargin}{-2cm}

% paragrafo
\addtolength{\parskip}{.2cm}

% Set up the header and footer
\pagestyle{fancy}
\lhead{\hmwkClass: \hmwkTitle} % Top center header
\rhead{\firstxmark} % Top right header
\lfoot{Maria Koren} % Bottom left footer
\cfoot{} % Bottom center footer
\rfoot{} % Bottom right footer
\renewcommand{\headrulewidth}{1pt}
\renewcommand{\footrulewidth}{1pt}

    
\newcommand{\hmwkTitle}{Nonogram za pomocą algorytmów genetycznych} 
\newcommand{\hmwkDueDate}{\today} 
\newcommand{\hmwkClass}{Inteligencja Obliczeniowa}
\newcommand{\hmwkAuthorName}{Maria Koren}

% trabalho 
\begin{document}
% capa
\begin{titlepage}
    \vfill
	\begin{center}
	\hspace*{-1cm}
	\vspace*{0.5cm}
    
\includegraphics[scale=0.55]{images/loga.png}\\
	\textbf{Uniwersytet Gdański \\ [0.05cm]Wydział Matematyki, Fizyki i Informatyki \\ [0.05cm] Instytut Informatyki}

	\vspace{0.6cm}
	\vspace{4cm}
	{\huge \textbf{\hmwkTitle}}\vspace{8mm}
	
	{\large \textbf{\hmwkAuthorName}}\\[3cm]
	
		\hspace{.45\textwidth} %posiciona a minipage
	   \begin{minipage}{.5\textwidth}

	  \end{minipage}
	  \vfill
	%\vspace{2cm}
	
	\textbf{Gdańsk}
	
	\textbf{\hmwkDueDate}
	\end{center}
	
\end{titlepage}

\newpage
\setcounter{secnumdepth}{5}
\tableofcontents
\newpage

\section{Wstęp}

W tym projekcie jest przedstawiony problem nonogram.

Nonogram (malowanie liczbami) to łamigłówka polegająca na zaczernianiu pól diagramu. Zaczernione pola utworzą rysunek. To, które pola trzeba zaczernić, wskazują liczby obok diagramu. Liczby z lewej strony każdego wiersza określają, ile grup czarnych pól jest w danym rzędzie i ile czarnych pól jest w każdej grupie.Analogicznie jest z liczbami u góry diagramu. 

Powyższy problem rozwiązano za pomocą algorytmu genetycznego (użyto 2 fukcje fitness) dla trzech rozmiarów: 5x5, 10x10, 15x15. 

Dodatkowo napisana mała aplikacja dla nonogramu 5x5.

\newpage
\section{Algorytm}

Dla powyższego problemu wybrano geny zero-jedynkowe, gdzie 1 oznacza zamalowaną klatkę, 0 pustą.

Użyto funkcji pomocniczej dla funkcji fitness:

\begin{lstlisting}[language=Python]
def fitness_sequence(sequence, clues):
        fitness = 0
        seq_id = 0
        
        for clue in clues:        
            #  Skip zeros in sequence
            while seq_id < len(sequence) and sequence[seq_id] == 0:
                seq_id += 1
                
            if seq_id == len(sequence):
                break
            
            # Count ones and check it matches to clue
            block_len = 0
            while seq_id < len(sequence) and sequence[seq_id] == 1:
                block_len += 1
                seq_id += 1   
            if block_len == clue:
                # +1 point for each correct block
                fitness += 1
        
        if(seq_id < len(sequence)):
            # -5 point for each additional block
            fitness -= np.sum(sequence[seq_id:])*5
        
        return fitness
\end{lstlisting}

Ta funkcja zlicza poprawne wartości i daje +1 punkt za każdą poprawną zgadniętą wartość i karę -5 za każdą niepoprawną.

Główna funkcja \textit{fitness}:
\begin{lstlisting}[language=Python]
 def fitness_func(model, solution, solution_idx):
        # Splits nonogram for columns and rows
        col_clues = np.array(nonogram[0], dtype=object)
        row_clues = np.array(nonogram[1], dtype=object)
        
        # Splits soluton to arrays of nonogram size
        grid = solution.reshape((len(row_clues), len(col_clues)))
        score = 0
        
        #  Checks each column solution
        for i in range(len(col_clues)):
            col_score = fitness_sequence(grid[:, i], col_clues[i])
            score += col_score
            
        # Checks each row solution
        for j in range(len(row_clues)):
            row_score = fitness_sequence(grid[j, :], row_clues[j])
            score += row_score
        return score
\end{lstlisting}

Główna funkcja fitness wywołuje pomicniczą dla każdej kolumny i każdego wiersza i zlicza sumarną wartość fitness (score).

Zrobiono również program z inną funkcją \textit{fitness}, która od poprzedniej różni się karą: tym razem za niepoprawne zgadywanie kara wynosi -1 pkt.

\newpage
\section{Porówanie}

Żeby porównać wyniki funkcji, użyto po 3 przykłady na każdy rozmiar. Wyniki są reprezentowane w tabeli \ref{tab:nonogram_results}.

\begin{table}[h]
    \centering
    \begin{tabular}{|c|c|c|c|c|}
        \hline
        Rozmiar & Algorytm & Statystyka & Wartości & Średnia \\ \hline
        \multirow{6}{*}{5x5} 
        & \multirow{2}{*}{nonogram\_ga} & Accuracy & 100.0, 100.0, 83.33 & 94.44 \\ \cline{3-5}
        & & Time & 1.91, 1.60, 3.47 & 2.33 \\ \cline{2-5}
        & \multirow{2}{*}{nonogram\_ga2} & Accuracy & 100.0, 100.0, 83.33 & 94.44 \\ \cline{3-5}
        & & Time & 10.33, 1.82, 3.18 & 5.11 \\ \hline
        \multirow{6}{*}{10x10} 
        & \multirow{2}{*}{nonogram\_ga} & Accuracy & 70.37, 82.98, 92.68 & 82.01 \\ \cline{3-5}
        & & Time & 23.98, 24.58, 26.17 & 24.91 \\ \cline{2-5}
        & \multirow{2}{*}{nonogram\_ga2} & Accuracy & 77.78, 87.23, 85.37 & 83.46 \\ \cline{3-5}
        & & Time & 27.38, 32.27, 33.64 & 31.10 \\ \hline
        \multirow{6}{*}{15x15} 
        & \multirow{2}{*}{nonogram\_ga} & Accuracy & 55.84, 70.59, 43.1 & 56.51 \\ \cline{3-5}
        & & Time & 49.84, 102.65, 50.30 & 67.60 \\ \cline{2-5}
        & \multirow{2}{*}{nonogram\_ga2} & Accuracy & 66.23, 78.82, 50.0 & 65.02 \\ \cline{3-5}
        & & Time & 55.45, 111.30, 52.62 & 73.12 \\ \hline
    \end{tabular}
    \caption{Wyniki dla różnych algorytmów i rozmiarów nonogramów}
    \label{tab:nonogram_results}
\end{table}

Na podstawie przedstawionych wyników można stwierdzić, że algorytm 1 zawsze działa szybciej niż algorytm 2, i daje nie gorsze wyniki.

\newpage
\section{Krzywy uczenia się}

Dla fittnes 1:

\begin{figure}[h]
    \centering
    \begin{minipage}{0.3\textwidth}
        \centering
        \includegraphics[width=\linewidth]{images/nonogram_ga_1_1.png}
        \caption{5x5, przykład 1}
    \end{minipage}\hfill
    \begin{minipage}{0.3\textwidth}
        \centering
        \includegraphics[width=\linewidth]{images/nonogram_ga_1_2.png}
        \caption{5x5, przykład 2}
    \end{minipage}\hfill
    \begin{minipage}{0.3\textwidth}
        \centering
        \includegraphics[width=\linewidth]{images/nonogram_ga_1_3.png}
        \caption{5x5, przykład 3}
    \end{minipage}

    \vspace{0.5cm} % odstęp między wierszami

    \begin{minipage}{0.3\textwidth}
        \centering
        \includegraphics[width=\linewidth]{images/nonogram_ga_2_1.png}
        \caption{10x10, przykład 1}
    \end{minipage}\hfill
    \begin{minipage}{0.3\textwidth}
        \centering
        \includegraphics[width=\linewidth]{images/nonogram_ga_2_2.png}
        \caption{10x10, przykład 2}
    \end{minipage}\hfill
    \begin{minipage}{0.3\textwidth}
        \centering
        \includegraphics[width=\linewidth]{images/nonogram_ga_2_3.png}
        \caption{10x10, przykład 3}
    \end{minipage}

    \vspace{0.5cm} % odstęp między wierszami

    \begin{minipage}{0.3\textwidth}
        \centering
        \includegraphics[width=\linewidth]{images/nonogram_ga_3_1.png}
        \caption{15x15, przykład 1}
    \end{minipage}\hfill
    \begin{minipage}{0.3\textwidth}
        \centering
        \includegraphics[width=\linewidth]{images/nonogram_ga_3_2.png}
        \caption{15x15, przykład 2}
    \end{minipage}\hfill
    \begin{minipage}{0.3\textwidth}
        \centering
        \includegraphics[width=\linewidth]{images/nonogram_ga_3_3.png}
        \caption{15x15, przykład 3}
    \end{minipage}
\end{figure}

\newpage

Dla fitness 2:
\begin{figure}[h]
    \centering
    \begin{minipage}{0.3\textwidth}
        \centering
        \includegraphics[width=\linewidth]{images/nonogram_ga2_1_1.png}
        \caption{5x5, przykład 1}
    \end{minipage}\hfill
    \begin{minipage}{0.3\textwidth}
        \centering
        \includegraphics[width=\linewidth]{images/nonogram_ga2_1_2.png}
        \caption{5x5, przykład 2}
    \end{minipage}\hfill
    \begin{minipage}{0.3\textwidth}
        \centering
        \includegraphics[width=\linewidth]{images/nonogram_ga2_1_3.png}
        \caption{5x5, przykład 3}
    \end{minipage}

    \vspace{0.5cm} % odstęp między wierszami

    \begin{minipage}{0.3\textwidth}
        \centering
        \includegraphics[width=\linewidth]{images/nonogram_ga2_1_2.png}
        \caption{10x10, przykład 1}
    \end{minipage}\hfill
    \begin{minipage}{0.3\textwidth}
        \centering
        \includegraphics[width=\linewidth]{images/nonogram_ga2_2_2.png}
        \caption{10x10, przykład 2}
    \end{minipage}\hfill
    \begin{minipage}{0.3\textwidth}
        \centering
        \includegraphics[width=\linewidth]{images/nonogram_ga2_2_3.png}
        \caption{10x10, przykład 3}
    \end{minipage}

    \vspace{0.5cm} % odstęp między wierszami

    \begin{minipage}{0.3\textwidth}
        \centering
        \includegraphics[width=\linewidth]{images/nonogram_ga2_3_1.png}
        \caption{15x15, przykład 1}
    \end{minipage}\hfill
    \begin{minipage}{0.3\textwidth}
        \centering
        \includegraphics[width=\linewidth]{images/nonogram_ga2_3_2.png}
        \caption{15x15, przykład 2}
    \end{minipage}\hfill
    \begin{minipage}{0.3\textwidth}
        \centering
        \includegraphics[width=\linewidth]{images/nonogram_ga2_3_3.png}
        \caption{15x15, przykład 3}
    \end{minipage}
\end{figure}

\newpage
\section{Aplikacja}

Przedstawione zostały wyniki dla 3 nonogramów każdego rozmiaru, a przetestowano było więcej. Zauważono, że dla nonogramu 5x5 wyniki często wychodzą 100\%. Oraz wyniki są otrzymywane dość szybko. Dlatego postanowiono napisać aplikacje właśnie dla tego rozmiaru.

Dla aplikcji wzięto algorytm z pierwszą funkcją fitness, ponieważ dokładność jest ta sama, natomiast jest szybsza, co w sytuacji aplikacji ma znaczenie.

W celu uruchamienia aplikcji trzeba wejść do folderu server i użyć polecenia:
\begin{lstlisting}[language=Bash]
    python main.py
\end{lstlisting}

Zatem wejśc w przeglądarkę pod strone \textit{localhost:5000} i będzie widoczna strona:

\begin{figure}[h]
    \centering
    \includegraphics[width=0.3\textwidth]{images/initsite.png}
\end{figure}

Podając liczby, opisujący chcące zamalowanie i naciskając przycisk \textbf{submit} otrzymujemy niżej wyniki. Przykład poniżej:

\begin{figure}[h]
    \centering
    \includegraphics[width=0.3\textwidth]{images/resultsite.png}
\end{figure}

\newpage
\section{Podsumowanie}
Zostały napisane 2 programy używające algorytmy genetyczne dla rozwiązania problemu nonogram. Porównując algorytmy, wyszło, że jeden algorytm działa szybciej, ale daje gorsze wyniki. Drugi na odwrót działa wolniej, ale daje lepsze wyniki.

Ponadto dla problemu w rozmiarze 5x5 napisano aplikacje, która używa ten algorytm rozwiązywania nonogramu.
\end{document}